% GENERATED FILE. Edit canonical lessons, then run npm run generate:books.
% canonical-source-hash: fnv1a64:d9a3d968e1108b88
% canonical-lessons: KA-C379-sammatisu, KA-C379-bhinnabhipraya, KA-C379-ommata, KA-C379-virodhisu, KA-C379-paraspara

\chapter{Agreeing and Disagreeing}
\label{ch:ka-agreeing-and-disagreeing}
\hlchaptermodality{379}

\begin{chapteropening}
\textbf{By the end of this chapter:} \emph{I can say to consent, a difference of opinion, a consensus, to oppose and mutual.}

Complete the last lesson of chapter 379: I can say to consent, a difference of opinion, a consensus, to oppose and mutual.
\end{chapteropening}

\section[\texorpdfstring{\kn{ಸಮ್ಮತಿಸು} (sammatisu)}{ಸಮ್ಮತಿಸು (sammatisu)}]{\kn{ಸಮ್ಮತಿಸು} (sammatisu) — to consent}

\label{lesson:KA-C379-sammatisu}



\begin{quote}
Before the new one: say the Kannada for bias, then the Kannada for to advocate.
\end{quote}

\subsection*{You'll want to know: \kn{ಸಮ್ಮತಿಸು}}
\textbf{\kn{ಸಮ್ಮತಿಸು}} — \emph{sammatisu} — \textquotedblleft{}to consent\textquotedblright{}.

\begin{quote}
\kn{ಅಪ್ಪ} \kn{ನನ್ನ} \kn{ಪ್ರವಾಸಕ್ಕೆ} \kn{ಸಮ್ಮತಿಸಿದರು}.
\end{quote}

\begin{quote}
\emph{appa nanna pravāsakke sammatisidaru.}
\end{quote}

\begin{quote}
Father agreed to my trip.
\end{quote}

Past (I …): \kn{ಸಮ್ಮತಿಸಿದೆ}, \emph{sammatiside}.

What you consent to takes -\kn{ಗೆ}. \kn{ಸಮ್ಮತಿ}, consent.

\begin{cousinweb}
\textbf{\kn{ಸಮ್ಮತಿಸು}} (\emph{sammatisu}) is built on Sanskrit \emph{sammati}, \textquotedblleft{}agreement, approval\textquotedblright{} (Monier-Williams), with the Kannada verb ending -\emph{isu}.
\end{cousinweb}

\begin{grammarlens}[title={What you've built}]
One more word for giving reasons and opinions.
\end{grammarlens}

\subsection*{Your turn}
\begin{itemize}
\raggedright
  \item \emph{Say it:} \emph{sammatisu}
  \item \emph{Say it:} \emph{sammatisu}, once more
  \item \emph{Say it:} \emph{appa nanna pravāsakke sammatisidaru.} — Father agreed to my trip.
  \item \emph{Recall:} say the Kannada for bias, then the Kannada for to advocate, then say \emph{sammatisu} again
\end{itemize}

\begin{tcolorbox}[breakable,colback=teal!4,colframe=teal!35!black,title={Before you move on}]
What does \kn{ಸಮ್ಮತಿಸು} mean? (\textquotedblleft{}To consent\textquotedblright{}.) Say it once more.
\end{tcolorbox}
\section[\texorpdfstring{\kn{ಭಿನ್ನಾಭಿಪ್ರಾಯ} (bhinnābhiprāya)}{ಭಿನ್ನಾಭಿಪ್ರಾಯ (bhinnābhiprāya)}]{\kn{ಭಿನ್ನಾಭಿಪ್ರಾಯ} (bhinnābhiprāya) — a difference of opinion}

\label{lesson:KA-C379-bhinnabhipraya}



\begin{quote}
Before the new one: say the Kannada for to advocate, then the Kannada for to consent.
\end{quote}

\subsection*{You'll want to know: \kn{ಭಿನ್ನಾಭಿಪ್ರಾಯ}}
\textbf{\kn{ಭಿನ್ನಾಭಿಪ್ರಾಯ}} — \emph{bhinnābhiprāya} — \textquotedblleft{}a difference of opinion\textquotedblright{}.

\begin{quote}
\kn{ಅಣ್ಣ} \kn{ತಮ್ಮನ} \kn{ನಡುವೆ} \kn{ಸಣ್ಣ} \kn{ಭಿನ್ನಾಭಿಪ್ರಾಯ} \kn{ಇತ್ತು}.
\end{quote}

\begin{quote}
\emph{aṇṇa tammana naḍuve saṇṇa bhinnābhiprāya ittu.}
\end{quote}

\begin{quote}
There was a small difference of opinion between the brothers.
\end{quote}

\kn{ಭಿನ್ನ}, different, + \kn{ಅಭಿಪ್ರಾಯ}, opinion; the two vowels merge into \kn{ಆ}.

\begin{cousinweb}
\textbf{\kn{ಭಿನ್ನಾಭಿಪ್ರಾಯ}} (\emph{bhinnābhiprāya}) is built from Sanskrit \emph{bhinna}, \textquotedblleft{}split, different\textquotedblright{} (CDIAL 9498), and \emph{abhiprāya}, \textquotedblleft{}aim, opinion\textquotedblright{} (Monier-Williams).
\end{cousinweb}

\begin{grammarlens}[title={What you've built}]
One more word for giving reasons and opinions.
\end{grammarlens}

\subsection*{Your turn}
\begin{itemize}
\raggedright
  \item \emph{Say it:} \emph{bhinnābhiprāya}
  \item \emph{Say it:} \emph{bhinnābhiprāya}, once more
  \item \emph{Say it:} \emph{aṇṇa tammana naḍuve saṇṇa bhinnābhiprāya ittu.} — There was a small difference of opinion between the brothers.
  \item \emph{Recall:} say the Kannada for to advocate, then the Kannada for to consent, then say \emph{bhinnābhiprāya} again
\end{itemize}

\begin{tcolorbox}[breakable,colback=teal!4,colframe=teal!35!black,title={Before you move on}]
What does \kn{ಭಿನ್ನಾಭಿಪ್ರಾಯ} mean? (\textquotedblleft{}A difference of opinion\textquotedblright{}.) Say it once more.
\end{tcolorbox}
\section[\texorpdfstring{\kn{ಒಮ್ಮತ} (ommata)}{ಒಮ್ಮತ (ommata)}]{\kn{ಒಮ್ಮತ} (ommata) — a consensus}

\label{lesson:KA-C379-ommata}



\begin{quote}
Before the new one: say the Kannada for to consent, then the Kannada for a difference of opinion.
\end{quote}

\subsection*{You'll want to know: \kn{ಒಮ್ಮತ}}
\textbf{\kn{ಒಮ್ಮತ}} — \emph{ommata} — \textquotedblleft{}a consensus\textquotedblright{}.

\begin{quote}
\kn{ಕೊನೆಗೆ} \kn{ಎಲ್ಲರಲ್ಲೂ} \kn{ಒಮ್ಮತ} \kn{ಮೂಡಿತು}.
\end{quote}

\begin{quote}
\emph{konege ellarallū ommata mūḍitu.}
\end{quote}

\begin{quote}
At last a consensus formed among everyone.
\end{quote}

\kn{ಒಮ್ಮತದಿಂದ}, unanimously.

\begin{cousinweb}
\textbf{\kn{ಒಮ್ಮತ}} (\emph{ommata}) is built from the native base for \textquotedblleft{}one\textquotedblright{}, \emph{or-} / \emph{om-} (DEDR 990), joined to Sanskrit \emph{mata}, \textquotedblleft{}an opinion\textquotedblright{} (Monier-Williams): one opinion.
\end{cousinweb}

\begin{grammarlens}[title={What you've built}]
One more word for giving reasons and opinions.
\end{grammarlens}

\subsection*{Your turn}
\begin{itemize}
\raggedright
  \item \emph{Say it:} \emph{ommata}
  \item \emph{Say it:} \emph{ommata}, once more
  \item \emph{Say it:} \emph{konege ellarallū ommata mūḍitu.} — At last a consensus formed among everyone.
  \item \emph{Recall:} say the Kannada for to consent, then the Kannada for a difference of opinion, then say \emph{ommata} again
\end{itemize}

\begin{tcolorbox}[breakable,colback=teal!4,colframe=teal!35!black,title={Before you move on}]
What does \kn{ಒಮ್ಮತ} mean? (\textquotedblleft{}A consensus\textquotedblright{}.) Say it once more.
\end{tcolorbox}
\section[\texorpdfstring{\kn{ವಿರೋಧಿಸು} (virōdhisu)}{ವಿರೋಧಿಸು (virōdhisu)}]{\kn{ವಿರೋಧಿಸು} (virōdhisu) — to oppose}

\label{lesson:KA-C379-virodhisu}



\begin{quote}
Before the new one: say the Kannada for a difference of opinion, then the Kannada for a consensus.
\end{quote}

\subsection*{You'll want to know: \kn{ವಿರೋಧಿಸು}}
\textbf{\kn{ವಿರೋಧಿಸು}} — \emph{virōdhisu} — \textquotedblleft{}to oppose\textquotedblright{}.

\begin{quote}
\kn{ಹಲವರು} \kn{ಹೊಸ} \kn{ತೆರಿಗೆಯನ್ನು} \kn{ವಿರೋಧಿಸಿದರು}.
\end{quote}

\begin{quote}
\emph{halavaru hosa terigeyannu virōdhisidaru.}
\end{quote}

\begin{quote}
Many people opposed the new tax.
\end{quote}

Past (I …): \kn{ವಿರೋಧಿಸಿದೆ}, \emph{virōdhiside}.

What you oppose takes -\kn{ಅನ್ನು}. \kn{ವಿರೋಧ}, opposition.

\begin{cousinweb}
\textbf{\kn{ವಿರೋಧಿಸು}} (\emph{virōdhisu}) is built on Sanskrit \emph{virodha}, \textquotedblleft{}opposition, hostility\textquotedblright{} (Monier-Williams), with the Kannada verb ending -\emph{isu}.
\end{cousinweb}

\begin{grammarlens}[title={What you've built}]
One more word for giving reasons and opinions.
\end{grammarlens}

\subsection*{Your turn}
\begin{itemize}
\raggedright
  \item \emph{Say it:} \emph{virōdhisu}
  \item \emph{Say it:} \emph{virōdhisu}, once more
  \item \emph{Say it:} \emph{halavaru hosa terigeyannu virōdhisidaru.} — Many people opposed the new tax.
  \item \emph{Recall:} say the Kannada for a difference of opinion, then the Kannada for a consensus, then say \emph{virōdhisu} again
\end{itemize}

\begin{tcolorbox}[breakable,colback=teal!4,colframe=teal!35!black,title={Before you move on}]
What does \kn{ವಿರೋಧಿಸು} mean? (\textquotedblleft{}To oppose\textquotedblright{}.) Say it once more.
\end{tcolorbox}
\section[\texorpdfstring{\kn{ಪರಸ್ಪರ} (paraspara)}{ಪರಸ್ಪರ (paraspara)}]{\kn{ಪರಸ್ಪರ} (paraspara) — mutual; each other}

\label{lesson:KA-C379-paraspara}



\begin{quote}
Before the new one: say the Kannada for a consensus, then the Kannada for to oppose.
\end{quote}

\subsection*{You'll want to know: \kn{ಪರಸ್ಪರ}}
\textbf{\kn{ಪರಸ್ಪರ}} — \emph{paraspara} — \textquotedblleft{}mutual; each other\textquotedblright{}.

\begin{quote}
\kn{ಅವರು} \kn{ಪರಸ್ಪರ} \kn{ಸಹಾಯ} \kn{ಮಾಡುತ್ತಾರೆ}.
\end{quote}

\begin{quote}
\emph{avaru paraspara sahāya māḍuttāre.}
\end{quote}

\begin{quote}
They help each other.
\end{quote}

Before a verb it means each other: \kn{ಪರಸ್ಪರ} \kn{ಮಾತಾಡು}, to talk to each other.

\begin{cousinweb}
\textbf{\kn{ಪರಸ್ಪರ}} (\emph{paraspara}) is borrowed from Sanskrit \emph{paraspara}, \textquotedblleft{}mutual, each other\textquotedblright{} (Monier-Williams), a doubling of \emph{para}, \textquotedblleft{}other\textquotedblright{}.
\end{cousinweb}

\begin{grammarlens}[title={What you've built}]
One more word for giving reasons and opinions.
\end{grammarlens}

\subsection*{Your turn}
\begin{itemize}
\raggedright
  \item \emph{Say it:} \emph{paraspara}
  \item \emph{Say it:} \emph{paraspara}, once more
  \item \emph{Say it:} \emph{avaru paraspara sahāya māḍuttāre.} — They help each other.
  \item \emph{Recall:} say the Kannada for a consensus, then the Kannada for to oppose, then say \emph{paraspara} again
\end{itemize}

\begin{tcolorbox}[breakable,colback=teal!4,colframe=teal!35!black,title={Before you move on}]
What does \kn{ಪರಸ್ಪರ} mean? (\textquotedblleft{}Mutual; each other\textquotedblright{}.) Say it once more.
\end{tcolorbox}
